\documentclass[11pt,a4paper]{article}
\usepackage[utf8]{inputenc}
\usepackage[T1]{fontenc}
\usepackage[spanish]{babel}
\usepackage{amsmath, amsthm, amssymb}
\usepackage{xcolor}
 \usepackage{tikz}
% --- Tipografía sans-serif ---
\usepackage{lmodern}
\renewcommand{\familydefault}{\sfdefault}
 
% Gráficos 2D y 3D
\usepackage{tikz}
\usepackage{pgfplots}
\pgfplotsset{compat=1.18, colormap/viridis}
 
% Cuadros personalizados, sombras y entornos numerados
\usepackage[many]{tcolorbox}
\tcbuselibrary{skins, breakable, theorems}

\usetikzlibrary{babel}
\usepackage{etoc} 

\newcommand{\portadaseccion}[1]{
    \clearpage
    \thispagestyle{empty}
    \vspace*{2cm}
    \section{#1} 
    \vspace{1.5cm} 
    \etocsettocstyle{\Large\bfseries Contenido del tema\par\vspace{0.5cm}}{}
    \localtableofcontents
    \clearpage 
}
 
% --- REFERENCIAS Y CITAS ENLAZADAS ---
\usepackage{hyperref}
\usepackage[capitalise, spanish]{cleveref}
\hypersetup{
    colorlinks   = true,
    linkcolor    = blue!75!black,
    citecolor    = violet!60!black,
    urlcolor     = cyan!70!black,
    pdftitle     = {Apuntes de Topología Algebraica},
    pdfborder    = {0 0 0},
    bookmarksnumbered = true,
}
 
% Nombres para que \cref reconozca las cajas personalizadas
\crefname{tcb@cnt@definicion}{definición}{definiciones}
\crefname{tcb@cnt@observacion}{observación}{observaciones}
\crefname{tcb@cnt@nota}{nota}{notas}
\crefname{tcb@cnt@teorema}{teorema}{teoremas}
\crefname{tcb@cnt@proposicion}{proposición}{proposiciones}
\Crefname{tcb@cnt@definicion}{Definición}{Definiciones}
\Crefname{tcb@cnt@observacion}{Observación}{Observaciones}
\Crefname{tcb@cnt@nota}{Nota}{Notas}
\Crefname{tcb@cnt@teorema}{Teorema}{Teoremas}
\Crefname{tcb@cnt@proposicion}{Proposición}{Proposiciones}
 
% --- Fondo y texto ---
\pagecolor{gray!5}
\color{black}
 
% --- Línea de degradado ("scan line") bajo cada \section (CORREGIDA) ---
\usepackage{titlesec}
\newcommand{\neonrule}[1]{%
    \par\noindent % Fuerza el origen al margen izquierdo
    \begin{tikzpicture}[remember picture, overlay]
    % 3 Bloques izquierdos
    \fill[#1] (0,0.3) rectangle (0.2,0.45);
    \fill[#1] (0.25,0.3) rectangle (0.45,0.45);
    \fill[#1] (0.5,0.3) rectangle (0.7,0.45);
    
    % Línea central 
    \draw[#1, line width=1pt] (1.2,0.375) -- (\dimexpr\linewidth-0.4cm\relax,0.375);
    
    % Bloque derecho
    \fill[#1] (\dimexpr\linewidth-0.5cm\relax, 0.3) rectangle (\linewidth,0.45);
\end{tikzpicture}%
}
\titleformat{\section}
  {\normalfont\Large\bfseries}{\thesection}{1em}{}
  [\neonrule{violet}] % Color lila para la línea de la sección
\titlespacing*{\section}{0pt}{3.5ex plus 1ex minus .2ex}{2.3ex plus .2ex}
 
% --- Sombra con brillo de color + título en degradado ---
\tcbset{
    futuristicbase/.style={
        enhanced, breakable,
        boxrule=0.5pt,
        arc=2pt,
        fonttitle=\bfseries
    }
}
 
% --- Cuadros personalizados numerados por sección (Paleta Azul y Lila) ---
 
\newtcbtheorem[number within=section]{teorema}{Teorema}{
    futuristicbase,
    colback=blue!3!white,
    colframe=blue!80!black,
    fuzzy shadow={1.2mm}{-1.2mm}{0mm}{0.3mm}{blue!50!black,opacity=0.5},
    title style={left color=blue!90!black, right color=blue!60!black}
}{thm}

\newtcbtheorem[number within=section]{definicion}{Definición}{
    futuristicbase,
    colback=violet!4!white,
    colframe=violet!80!black,
    fuzzy shadow={1.2mm}{-1.2mm}{0mm}{0.3mm}{violet!50!black,opacity=0.5},
    title style={left color=violet!80!black, right color=violet!50!black}
}{def}

\newtcbtheorem[number within=section]{proposicion}{Proposición}{
    futuristicbase,
    colback=blue!2!violet!3!white,
    colframe=blue!50!violet!50!black,
    fuzzy shadow={1.2mm}{-1.2mm}{0mm}{0.3mm}{blue!40!violet!40!black,opacity=0.5},
    title style={left color=blue!60!violet!60!black, right color=blue!40!violet!40!black}
}{prop}

\newtcbtheorem[number within=section]{observacion}{Observación}{
    futuristicbase,
    colback=purple!5!white,
    colframe=purple!50!black,
    boxrule=0.4pt,
    coltitle=black,
    fuzzy shadow={1.2mm}{-1.2mm}{0mm}{0.3mm}{purple!30!black,opacity=0.5},
    title style={left color=purple!20!white, right color=purple!5!white}
}{obs}

\newtcbtheorem[number within=section]{nota}{Nota}{
    futuristicbase,
    colback=magenta!5!blue!5!white,
    colframe=magenta!60!blue!60!black,
    boxrule=0.4pt,
    coltitle=white,
    fuzzy shadow={1.2mm}{-1.2mm}{0mm}{0.3mm}{magenta!40!blue!40!black,opacity=0.5},
    title style={left color=magenta!80!blue!80!black, right color=magenta!60!blue!60!black}
}{not}
 
% --- Cabecera/pie de página ---
\usepackage{fancyhdr}
\setlength{\headheight}{12.5pt}
\pagestyle{fancy}
\fancyhf{}
\renewcommand{\headrulewidth}{0.4pt}
\renewcommand{\footrulewidth}{0.4pt}
\fancyhead[L]{\small\textcolor{gray}{\leftmark}}
\fancyhead[R]{\small\textcolor{gray}{Apuntes}}
\fancyfoot[C]{\small\textcolor{gray}{\thepage}}
 
